\documentclass[10pt,a4paper]{article}
\usepackage{fontspec}
\usepackage{polyglossia}
\setmainlanguage{russian}
\setmainfont{Liberation Serif}
\setmonofont{Liberation Mono}
\usepackage[left=1.6cm,right=1.6cm,top=1.3cm,bottom=1.3cm]{geometry}
\usepackage{amsmath,amssymb,xcolor,tabularx,array}
\setlength{\parindent}{0pt}\setlength{\parskip}{2pt}
\setlength{\emergencystretch}{3em}
\pagestyle{plain}


\newcommand{\usl}[1]{\par\vspace{7pt}\noindent\colorbox{black!7}{\parbox{\dimexpr\linewidth-2\fboxsep}{#1}}\par\vspace{2pt}}
\newcommand{\pole}[2]{\par\noindent\hangindent=1.2em\textbf{#1}\ #2\par}
\begin{document}
\noindent{\Large\bfseries Построй или докажи, что нельзя}\hfill{\small Кружок 7\,И · 29 сентября · разбор для преподавателей}
\vspace{3pt}\hrule\vspace{5pt}

\textbf{Идея.} «Можно»~--- пример с проверкой, «нельзя»~--- доказательство. Новое после логики: запрет даёт одна величина
на всю конструкцию (сумма, концы проводов, отметки, поражения), а не цепочка «рыцарь, значит, сосед лжец\,\ldots».

\vspace{3pt}
\noindent\small\begin{tabularx}{\linewidth}{@{}l l X@{}}
\textbf{1--3} & вход & 1~--- любой пример; 2~--- кажется «нельзя», а можно; 3~--- первый запрет одной величиной и где он ломается\\
\textbf{4--7} & середина & 4~--- чётность, «при каких $n$»; 5~--- шаг $+3$ и запрет углами; 6, 7~--- «найдите все»: примеры и запрет остального\\
\textbf{8$^*$--9$^*$} & верх & 8$^*$~--- запрет подсчётом; 9$^*$~--- углы и разбор случаев\\
\textbf{10$^*$} & вершина & собирает 4(а), 5(б) и «при каких $n$»\\
\end{tabularx}\normalsize

\vspace{3pt}
\textbf{Сдача:} пример с проверкой; «нельзя»~--- доказательство, не «не получилось». \textbf{Подсказка}~--- застрявшему, \textbf{Спросить}~--- при сдаче.
\textbf{Начало:} голосование «можно или нельзя» по задаче 2.
\textbf{Концовка у доски, без ответа:} доска $8\times8$ без двух противоположных угловых клеток~--- на доминошки $1\times2$? (Мостик к раскраскам.)

\usl{\textbf{1.} Разрежьте квадрат на \textbf{(а)}~7 квадратов; \textbf{(б)}~6 квадратов;
\textbf{(в)}~8 квадратов. Квадраты не обязаны быть одинаковыми.}
\pole{Ответ:}{всё можно.}
\pole{Решение.}{(а) Средние линии~--- 4 квадрата, один ещё на 4: $3+4=7$.\\ (б) Сетка $3\times3$, угол $2\times2$~--- один квадрат: $1+5=6$.\\ (в) Сетка $4\times4$, угол $3\times3$~--- один: $1+7=8$.}
\pole{Подсказка.}{Застрял на 6~--- «начни с 9 клеток». Сетка $k\times k$ со склеенным углом~--- $2k$ кусков.}
\pole{Спросить.}{Сильному: «а на 5?», без ответа (это 9$^*$).}
\usl{\textbf{2.} Могут ли и сумма, и произведение нескольких (хотя бы двух) натуральных чисел
равняться 99?}
\pole{Ответ:}{да, например 9, 11 и 79 единиц.}
\pole{Решение.}{$9\cdot11=99$, единицы произведение не меняют; $9+11+79=99$.}
\pole{Подсказка.}{Разложи 99 на множители, добей сумму единицами. Годится любое разложение (3, 33 и 63 единицы).}
\pole{Спросить.}{Сумму пересчитал?}
\usl{\textbf{3.} В ряд записаны \textbf{(а)}~10 чисел; \textbf{(б)}~9 чисел. Сумма любых двух соседних
чисел положительна. Может ли сумма всех чисел быть отрицательной?}
\pole{Ответ:}{(а) нет; (б) да.}
\pole{Решение.}{(а) Пять пар соседей, каждая $>0$.\\ (б) $-9,\,10,\,-9,\,\ldots,\,10,\,-9$ (пять $-9$, четыре 10): соседи дают 1, всего $-5$.}
\pole{Подсказка.}{В (б) одно число без пары~--- его и делаем очень отрицательным.}
\pole{Спросить.}{Почему довод (а) не проходит в (б). Сильному: при каких $n$?~--- при нечётных.}
{\small\pole{Источник:}{Малый мехмат, 7 класс, 2024.}}
\usl{\textbf{4.} Телефоны соединяют проводами, каждый провод соединяет два телефона,
и никакие два телефона не соединены дважды.\par
\textbf{(а)}~Можно ли соединить 15 телефонов так, чтобы каждый был соединён ровно с тремя другими?\par
\textbf{(б)}~При каких $n$ так можно соединить $n$ телефонов?}
\pole{Ответ:}{(а) нет; (б) при чётных $n\ge4$.}
\pole{Решение.}{(а) Концов $15\cdot3=45$, а их вдвое больше, чем проводов.\\ (б) Нечётные~--- как (а); $n=2$~--- меньше трёх соседей. Чётное $n\ge4$: по кругу, каждый соединён с двумя соседями и с противоположным (при $n\ge4$ он не сосед).}
\pole{Подсказка.}{Нет примера~--- «пусть все одинаковы: посади по кругу».}
\pole{Спросить.}{В (б) три части: нечётные, $n=2$, пример. Чаще теряют $n=2$.}
\usl{\textbf{5.} \textbf{(а)}~Разрежьте правильный треугольник на 6, на 7 и на 8 правильных треугольников
(не обязательно одинаковых).\par
\textbf{(б)}~Докажите, что его можно разрезать на $n$ правильных треугольников при любом $n\ge6$.\par
\textbf{(в)}~Докажите, что на 2 и на 3 правильных треугольника его разрезать нельзя.}
\pole{Ответ:}{(а), (б) можно; (в) нельзя.}
\pole{Решение.}{(а) Как в 1: 7~--- средние линии и один ещё на 4; 6 и 8~--- сетки из 9 и 16, у вершины склеить 4 и 9.\\ (б) Кусок на 4: $+3$; от 6, 7, 8~--- все $n\ge6$.\\ (в) В каждой вершине большого~--- кусок своим углом; кусок с двумя вершинами~--- весь треугольник. Значит, кусков $\ge3$, а при трёх (стороны $a,b,c$) каждую сторону большого закрывают два угловых: $a+b=b+c=c+a=1$, все по $\tfrac12$, в центре дыра.}
\pole{Подсказка.}{Застрял в (в)~--- «кто накрывает вершины?»}
\pole{Спросить.}{Почему хватает 6, 7, 8 (разные остатки при делении на 3). Сильному: полный ответ~--- 1, 4 и все $n\ge6$.}
\usl{\textbf{6.} На острове живут рыцари, которые всегда говорят правду, и лжецы,
которые всегда лгут. За круглым столом сидят островитяне, и каждый сказал:
\textit{«оба моих соседа~--- лжецы»}. Сколько рыцарей может сидеть за столом,
если всего островитян \textbf{(а)}~10; \textbf{(б)}~7? Найдите все варианты.}
\pole{Ответ:}{(а) 4 или 5; (б) 3.}
\pole{Решение.}{Рыцари не рядом; трёх лжецов подряд нет (средний сказал бы правду); все лжецы~--- тоже нет. Значит, между соседними рыцарями 1--2 лжеца: при $k$ рыцарях $2k\le n\le3k$.\\ (а) РЛРЛРЛЛРЛЛ ($k=4$), РЛРЛРЛРЛРЛ ($k=5$).\\ (б) РЛРЛРЛЛ ($k=3$).}
\pole{Подсказка.}{Застрял~--- «сколько лжецов может сидеть между двумя рыцарями?»}
\pole{Спросить.}{Оба ответа (а); почему не 3 и не 6. Сильному: при 12~--- три ответа (4, 5, 6).}
\usl{\textbf{7.} \textbf{(а)}~При каких $N$ и сумма, и произведение нескольких (хотя бы двух)
натуральных чисел могут равняться $N$?\par
\textbf{(б)}~Найдите все наборы из нескольких (хотя бы двух) \textit{различных} натуральных чисел,
у которых сумма равна произведению.}
\pole{Ответ:}{(а) при составных $N$; (б) только $1,2,3$.}
\pole{Решение.}{(а) $N=ab$, $a,b\ge2$: $a$, $b$ и $N-a-b\ge0$ единиц ($ab-a-b=(a-1)(b-1)-1$). Простое $N$: произведение $N$~--- только $N$ и единицы, сумма $>N$. $N=1$: сумма $\ge2$.\\ (б) $c$~--- наибольшее из $k$ чисел, $P$~--- произведение остальных. Сумма $<kc$ (остальные $<c$), произведение $cP$, значит, $P<k$; но $P\ge1\cdot2\cdots(k-1)$, так что $k\le3$. $k=2$: $(a-1)(b-1)=1$, $a=b=2$~--- не различны. $k=3$, $a<b<c$: $ab<3$~--- это $1,2$; $3+c=2c$, $c=3$.}
\pole{Подсказка.}{(а) «а если $N$ простое?» (б) пример~--- секунда, работа~--- «почему других нет»; «посмотри на наибольшее».}
\pole{Спросить.}{В (а) оба направления и $N=1$. В (б)~--- «почему нет четырёх чисел?»}
\usl{\textbf{8$^*$.} На кружке было 9 задач, и каждую из них рассказали преподавателю Аня, Боря и Вера,
каждый в свой момент. Аня отмечает у себя задачу, если рассказала её раньше Бори; Боря~--- если рассказал раньше Веры;
Вера~--- если рассказала раньше Ани. Может ли у каждого из троих оказаться отмечено \textbf{(а)}~по 6 задач; \textbf{(б)}~по 7 задач?}
\pole{Ответ:}{(а) да; (б) нет.}
\pole{Решение.}{(б) Три отметки за задачу~--- круг: А раньше Б, Б раньше В, В раньше А. Значит, $\le2$ на задачу, всего $\le18<21$.\\ (а) По три задачи в порядках АБВ, БВА, ВАБ: у каждого по 6.}
\pole{Подсказка.}{(б) подсказывает (а): 18~--- ровно по две отметки на задачу, а две дают только АБВ, БВА, ВАБ.}
\pole{Спросить.}{Почему не бывает трёх отметок; пересчёт отметок у всех троих.}
\pole{После.}{парадокс Кондорсе; кубики $(2,2,4,4,9,9)$, $(1,1,6,6,8,8)$, $(3,3,5,5,7,7)$~--- каждый бьёт следующий по кругу с вероятностью $\tfrac59$.}
{\small\pole{Источник:}{ЦПМ, «Хамовники», 7 класс, 2024, «Мнимые противоречия», №\,1; (б) добавлен.}}
\usl{\textbf{9$^*$.} Докажите, что квадрат нельзя разрезать на 5 квадратов.}
\pole{Ответ:}{доказательство.}
\pole{Решение.}{Сторона большого 1.\\ (1) Кусок в углу стоит в нём углом, стороны~--- по сторонам большого; кусок с двумя углами~--- весь квадрат. Значит, угловых четыре: $A$, $B$ сверху, $C$, $D$ снизу (стороны $a,b,c,d$); пятый~--- $E$.\\ (2) Сторону большого закрывают куски, лежащие на ней стороной: сумма их сторон~--- 1. $E$ лежит не более чем на одной стороне (на противоположных~--- его сторона $\ge1$, на соседних~--- занял бы угол).\\ (3) $E$ не у края: $a+b=c+d=a+c=b+d=1$, так что $a=d$, $b=c$. При $a\ne b$ противоположные угловые больше половины и перекрываются; при $a=b=\tfrac12$ места для $E$ нет.\\ (4) $E$ у верхней стороны (прочие так же): $a+e+b=1$, $a+c=c+d=b+d=1$, откуда $a=d$, $a+b=1$, $e=0$.}
\pole{Подсказка.}{«Кто накрывает углы?» (как 5(в)), потом «где пятый?»}
\pole{Спросить.}{Оба случая (3), (4) и почему $E$ не у двух сторон. Не принимать: «пятый маленький, как-то не выходит».}
\usl{\textbf{10$^*$.} В турнире не меньше двух участников, каждый сыграл с каждым по одному разу, ничьих не было. Назовём турнир
\textit{запутанным}, если для любых двух участников найдётся третий, который выиграл у обоих.\par
\textbf{(а)}~Докажите, что в запутанном турнире каждый участник проиграл хотя бы трём.\par
\textbf{(б)}~Докажите, что в запутанном турнире не меньше 7 участников.\par
\textbf{(в)}~Придумайте запутанный турнир с 7 участниками.\par
\textbf{(г)}~При каких $n$ существует запутанный турнир с $n$ участниками?}
\pole{Ответ:}{(г) при $n\ge7$.}
\pole{Решение.}{(а) $L$~--- обыгравшие $A$; $L$ не пусто (пара $A$ с кем угодно). Для $B\in L$ общий победитель $A$ и $B$ лежит в $L$ и обыграл $B$: в $L$ каждый проиграл кому-то из $L$. Один в $L$~--- некому; двое~--- победитель их партии никому в $L$ не проиграл. Значит, в $L$ хотя бы трое.\\ (б) Поражений $\ge3n$, а всего их $\tfrac{n(n-1)}2$ (одно на партию): $n\ge7$.\\ (в) $0,\ldots,6$ по кругу: $k$ бьёт $k+1$, $k+2$, $k+4$ (по модулю 7), остальным проигрывает. Сдвиг по кругу турнир не меняет, так что хватит пар $0,k$; их общие победители при $k=1,\ldots,6$: 6, 5, 6, 3, 3, 5.\\ (г) $n<7$~--- нет по (б). $n\to n+1$: копия $X^{\prime}$ играет со всеми как $X$, а $X$ бьёт $X^{\prime}$. Старые пары~--- прежние победители; $X,X^{\prime}$~--- любой обыгравший $X$; $X^{\prime},Y$~--- общий победитель $X$ и $Y$.}
\pole{Подсказка.}{(а) «посмотри на тех, кто обыграл $A$». (в) при $n=7$ в (б) равенство~--- каждый бьёт ровно троих, ищи по кругу. Ловушка: «бью следующих трёх» (у 0 и 3 нет общего победителя).}
\pole{Спросить.}{(а) почему в $L$ не двое; (в) все пары (чужой некруговой пример~--- все 21; годится и $k+3,k+5,k+6$); (г) обе новые пары.}
{\small\pole{Источник:}{проблема Шютте; минимум 7~--- Дж.~и~Э.~Секереш, 1965.}}

\vspace{8pt}\hrule\vspace{4pt}
\textbf{Резерв.}
\textit{Если решил всё:} $n$-значное число из цифр 1 и 2, делящееся на $2^n$, существует и единственно («Квант», 1971).
Приписываем цифру $d$ слева: $d\cdot10^n+2^nq=2^n(d\cdot5^n+q)$; $5^n$ нечётно, так что скобку делает чётной ровно одно из $d=1,2$~--- отсюда и существование, и единственность.
Труднее~--- нечётные цифры и $5^n$ (олимпиада США, 2003).
\textit{Другая фактура в середину:} числа $1,\ldots,12$ на рёбрах куба с равными суммами в вершинах~--- нельзя ($2\cdot78=156$ не делится на 8);
из $1,\ldots,13$ можно вычеркнуть только 3, 7 или 11, все три годятся (Турнир Ломоносова, 1978).
\end{document}
